Let $K$ be a field and let
\[
  \Omega=\{+,-,\times,/\}
\]
be the ordinary binary arithmetic signature.  Division is partial in the
affine chart: a denominator must be nonzero at the moment at which the
operation is performed.

An arithmetic expression is a planar binary tree.  When its internal vertices
form a single chain, its operations acquire a unique temporal order, and the
expression can be written in two extreme ways: the already-computed subtree may
occupy the first operand slot or the second.  This section fixes that
convention, proves the elementary facts on which every later chapter depends,
and isolates the one algebraic law that relates the two expansions of a common
expression.

\subsection{Chronological comb conventions}

\begin{definition}[Left- and right-expanded one-hole expressions]
\label{def:p0-left-right-combs}
Fix constants $c_1,\ldots,c_n\in K$ and operation labels
$\omega_1,\ldots,\omega_n\in\Omega$.  Starting from the hole $z$, define
recursively
\begin{equation}\label{eq:p0-comb-recursions}
  L_0[z]=z,
  \qquad
  L_k[z]=\omega_k(L_{k-1}[z],c_k),
\end{equation}
\begin{equation}\label{eq:p0-right-comb-recursions}
  R_0[z]=z,
  \qquad
  R_k[z]=\omega_k(c_k,R_{k-1}[z]).
\end{equation}
The expression $L_n$ is a \emph{left-expanded comb}; $R_n$ is a
\emph{right-expanded comb}.  The index $k$ is the evaluation time: the
innermost operation has index $1$ and is evaluated first.
\end{definition}

\begin{convention}[Scope of the words ``left'' and ``right'']
\label{conv:p0-left-right-convention}
Throughout this paper, ``left-expanded'' and ``right-expanded'' refer only to
the recursion \eqref{eq:p0-comb-recursions}--\eqref{eq:p0-right-comb-recursions}:
the active subtree occupies operand slot $(1)$ in the first case and operand
slot $(2)$ in the second.  The index is chronological.  Every later use of the
words must be reducible to \eqref{eq:p0-comb-recursions}--\eqref{eq:p0-right-comb-recursions};
the planar direction in which a tree happens to be drawn is not part of the
convention.
\end{convention}

Thus, in fully parenthesized form,
\[
  L_n[z]
  =\omega_n(\cdots\omega_2(\omega_1(z,c_1),c_2)\cdots,c_n),
\]
whereas
\[
  R_n[z]
  =\omega_n(c_n,\omega_{n-1}(c_{n-1},\cdots\omega_1(c_1,z)\cdots)).
\]
In the first expression the active subtree is the left child at every internal
vertex.  In the second it is the right child.

\begin{figure}[t]
\centering
\begin{tikzpicture}[
  op/.style={circle,draw=black!65,fill=blue!7,minimum size=7mm,inner sep=0pt},
  leaf/.style={rectangle,draw=black!45,fill=gray!5,minimum height=6mm,
              minimum width=8mm,inner sep=1pt},
  every node/.style={font=\small}
]
  \node[font=\bfseries] at (-3.2,2.65) {left-expanded};
  \node[op] (l3) at (-3.2,2.0) {$\omega_3$};
  \node[op] (l2) at (-4.0,1.15) {$\omega_2$};
  \node[leaf] (lc3) at (-2.35,1.15) {$c_3$};
  \node[op] (l1) at (-4.75,0.3) {$\omega_1$};
  \node[leaf] (lc2) at (-3.25,0.3) {$c_2$};
  \node[leaf] (lz) at (-5.35,-0.55) {$z$};
  \node[leaf] (lc1) at (-4.15,-0.55) {$c_1$};
  \draw (l3)--(l2); \draw (l3)--(lc3);
  \draw (l2)--(l1); \draw (l2)--(lc2);
  \draw (l1)--(lz); \draw (l1)--(lc1);
  \draw[-{Latex[length=2mm]},very thick,blue!60!black]
    (-5.75,-0.9) .. controls (-5.2,-1.25) and (-3.6,-1.25) .. (-2.8,-0.9)
    node[midway,below,font=\scriptsize] {evaluation follows the left spine outward};

  \node[font=\bfseries] at (3.2,2.65) {right-expanded};
  \node[op] (r3) at (3.2,2.0) {$\omega_3$};
  \node[leaf] (rc3) at (2.35,1.15) {$c_3$};
  \node[op] (r2) at (4.0,1.15) {$\omega_2$};
  \node[leaf] (rc2) at (3.25,0.3) {$c_2$};
  \node[op] (r1) at (4.75,0.3) {$\omega_1$};
  \node[leaf] (rc1) at (4.15,-0.55) {$c_1$};
  \node[leaf] (rz) at (5.35,-0.55) {$z$};
  \draw (r3)--(rc3); \draw (r3)--(r2);
  \draw (r2)--(rc2); \draw (r2)--(r1);
  \draw (r1)--(rc1); \draw (r1)--(rz);
  \draw[-{Latex[length=2mm]},very thick,orange!75!black]
    (5.75,-0.9) .. controls (5.2,-1.25) and (3.6,-1.25) .. (2.8,-0.9)
    node[midway,below,font=\scriptsize] {evaluation follows the right spine outward};
\end{tikzpicture}
\caption{The two pure combs under the chronological indexing convention.
Both have one chain of internal dependencies; the active subtree occupies
opposite operand slots.}
\label{fig:p0-two-combs}
\end{figure}

\begin{proposition}[Unique internal order]
\label{prop:p0-combs-unique-order}
Every finite left- or right-expanded comb has a unique legal order of internal
evaluation, namely the chronological order $1,2,\ldots,n$.
\end{proposition}

\begin{proof}
At stage $k$, the only internal child of the $k$th vertex is the subtree formed
at stage $k-1$; the other child is the leaf $c_k$.  Hence the internal vertices
form one dependency chain
\[
  v_1<v_2<\cdots<v_n.
\]
A finite chain has a unique linear extension.
\end{proof}

This proposition is deliberately weaker than the full sequential-tree
classification of Paper~I.  Here the combs are given by explicit grammars;
Paper~I proves intrinsically that every binary expression tree with a unique
internal evaluation order has one such spine after operand slots are recorded.

\subsection{One-hole sections and opposite operations}

For $\omega\in\Omega$ and $c\in K$, write
\[
  C^{(1)}_{\omega,c}[z]=\omega(z,c),
  \qquad
  C^{(2)}_{\omega,c}[z]=\omega(c,z).
\]
Then $L_n$ is a chronological word of pure slot-$(1)$ sections and $R_n$ a
word of pure slot-$(2)$ sections.

Define the opposite binary operation by
\[
  \omega^{\mathrm{op}}(u,v):=\omega(v,u).
\]

\begin{proposition}[Comb opposition]
\label{prop:p0-comb-opposition}
With chronological indexing, a right-expanded comb for operations
$\omega_1,\ldots,\omega_n$ is exactly a left-expanded comb for the opposite
operations:
\[
  R_k[z]
  =\omega_k^{\mathrm{op}}(R_{k-1}[z],c_k).
\]
Consequently, the two comb grammars are isomorphic after replacing every
operation by its opposite.  They agree literally for commutative operations.
\end{proposition}

\begin{proof}
By definition,
\[
  \omega_k(c_k,R_{k-1}[z])
  =\omega_k^{\mathrm{op}}(R_{k-1}[z],c_k).
\]
Apply this identity at every stage.  If $\omega_k$ is commutative, then
$\omega_k^{\mathrm{op}}=\omega_k$.
\end{proof}

The result is combinatorial, not an assertion that the two arithmetic maps are
always equal.  Their elementary one-hole maps are:

\begin{center}
\renewcommand{\arraystretch}{1.35}
\begin{tabular}{lll}
\toprule
Operation & active value in slot $(1)$ & active value in slot $(2)$ \\
\midrule
addition & $z+c$ & $c+z=z+c$ \\
subtraction & $z-c$ & $c-z$ \\
multiplication & $cz$ & $cz$ \\
division & $z/c$ & $c/z$ \\
\bottomrule
\end{tabular}
\end{center}

For $c\ne0$, every first-slot map is affine.  The second-slot subtraction is
also affine, but it has multiplier $-1$.  The second-slot division is not
affine: it has a finite pole at $z=0$ and sends that point to infinity under
projective continuation.

\begin{corollary}[Where the algebraic symmetry breaks]
\label{cor:p0-affine-projective-break}
If the operation alphabet is restricted to $+$ and $\times$, pure left and
pure right expansion have the same one-hole operator algebra.  With
subtraction they remain inside the affine group.  Once second-slot division
$c/z$ is admitted, the right-expanded language contains inversion and leaves
the affine stabilizer of infinity.
\end{corollary}

\begin{proof}
The first two assertions follow from the table.  For $c\ne0$,
\[
  \frac{c}{z}=D_c\circ J(z),
  \qquad D_c(z)=cz,\quad J(z)=\frac1z.
\]
The map $J$ exchanges $0$ and $\infty$, so it does not stabilize infinity.
\end{proof}

\subsection{The distributive law as an expansion rewrite}
\label{subsec:p0-distributivity}

Two different expansions may induce the same arithmetic operator.  The law used
here to relate two such expansions is distributivity, and it is worth isolating
because every later chapter compares expansions that share an operator or a
value.  Other ring identities (associativity, the unit laws) also relate
expansions with a common operator; what is special about distributivity is that
it moves an additive step across a multiplicative one, which is why it, and not
the others, produces the charge transport computed in
Example~\ref{ex:p0-distributive-charge-transport}.

\begin{definition}[Distributive expansion]
\label{def:p0-distributive-expansion}
Let $p,k\in K$.  The \emph{distributive expansion} of the two-step slot-$(1)$
word
\begin{equation}\label{eq:p0-distributive-expansion}
  z\;\longmapsto\;k(z+p)=(kz+kp)
\end{equation}
is the slot-$(1)$ word obtained by reading the right-hand side as a scaling
followed by a translation:
\[
  z\;\longmapsto\;kz\;\longmapsto\;kz+kp .
\]
The two words have the same length, the same operation alphabet, and the same
induced operator $z\mapsto kz+kp$ on $K$, but different constant data.  Writing
a two-step word chronologically as the pair (translation constant, scaling
constant), the first word is $(p,k)$ and the second is $(k,kp)$.
\end{definition}

\begin{proposition}[Distributivity is operator-preserving and expansion-changing]
\label{prop:p0-distributive-operator-identity}
For all $p,k\in K$ the induced operators of the two words in
\eqref{eq:p0-distributive-expansion} agree:
\begin{equation}\label{eq:p0-distributive-identity}
  \bigl(z\mapsto z+p\bigr)\ \text{followed by}\ \bigl(z\mapsto kz\bigr)
  \;=\;
  \bigl(z\mapsto kz\bigr)\ \text{followed by}\ \bigl(z\mapsto z+kp\bigr).
\end{equation}
The two words coincide as marked data precisely when $p=k=1$; for every other
pair they are distinct marked expansions.  In particular, if $p\ne0$ and
$k\ne1$, then the two expansions share their operator and differ in their
chronological constant data.
\end{proposition}

\begin{proof}
Both sides send $z$ to $kz+kp$, which is the distributive law in $K$.  The
chronological pairs are $(p,k)$ and $(k,kp)$, and the second differs from the
first unless $p=k$ and $k=kp$.  Since $k\ne0$ in the second equation and the
identity case $k=0$ is immediate in the first, the two conditions force
$p=1$ and then $k=1$.  Hence the words coincide as marked data only in the
trivial case.
\end{proof}

\begin{example}[Distributivity transports the additive charge]
\label{ex:p0-distributive-charge-transport}
Take $K=\mathbb R$, $p=1$, $k=e^{q}$ with $q\ne0$, and compare the two words
of Proposition~\ref{prop:p0-distributive-operator-identity}.  Both induce
$z\mapsto e^{q}z+e^{q}$.  Their accumulated charges, in the sense fixed in
Section~\ref{sec:p0-acs-tearing}, are
\[
  (A,M)=(1,q)
  \qquad\text{and}\qquad
  (A,M)=(e^{q},q).
\]
The multiplicative charge is the same because the two words differ only in
where the additive step is performed; the additive charge differs by exactly
the scale $e^{q}$ that separates the two positions.  In the notation of
Definition~\ref{def:acs}, the weighted integrals agree,
\[
  \int_{C_{\gamma}}e^{-M}\dd A
  =1\cdot e^{0}=1,
  \qquad
  \int_{C_{\delta}}e^{-M}\dd A
  =e^{q}\cdot e^{-q}=1,
\]
so that Proposition~\ref{prop:acs-evaluation} returns the common value
$e^{q}(x+1)$ along both charge paths.
\end{example}

Example~\ref{ex:p0-distributive-charge-transport} is the elementary core of the
later discussion of tearing.  Distributivity does not merely reorganize
symbols: it moves additive charge through multiplicative charge, and the moved
amount is exactly the scale accumulated in between.  An algebraic identity of
operators therefore becomes a statement about two different charge histories.

\begin{remark}[Where expansion stops]
\label{rem:p0-expansion-boundary}
Distributivity is a law of the ring $(\times,+)$ and not a universal law of the
whole signature $\Omega$.  Its analogues for the two remaining operations
behave differently, and the difference is instructive.

\begin{enumerate}
  \item For slot-$(2)$ subtraction,
  \[
    (c-z)k=kc-kz=(-k)z+kc .
  \]
  This is the distributive law again, and the transported constant is $kc$;
  the resulting one-hole map is still affine, but its multiplier is $-k$.
  Expansion therefore remains inside the affine sector.

  \item For slot-$(2)$ division,
  \[
    \frac{c}{z}\,k=\frac{kc}{z},
    \qquad z\ne0 .
  \]
  This is an identity of partial functions, not an expansion: the right-hand
  side is inversion followed by a scaling, and inversion is not an affine
  motion of the arithmetic value.  Reading it as a comb would require a
  generator outside the affine group.
\end{enumerate}

Consequently the expansion calculus of this section is an affine calculus, and
the projective step $c/z$ is exactly where it ends.  Section~\ref{sec:p0-ripples}
reads that step contravariantly instead, as a condition on outputs rather than
a motion of the value.
\end{remark}

\subsection{Expansion and mechanical computation}

The distributive law is what makes mechanical expansion possible at all.  A
symbolic engine that reduces a polynomial expression to a normal form does so
by distributing products over sums and collecting like terms; this is the
algebraic core of the mechanization of mathematics, whose systematic programme
was set out by Wu Wen-tsun \cite{Wu2000MathematicsMechanization}.  In that
tradition expansion is a means: the process is eliminated in favour of a
canonical object, and the mechanical content is transferred to an algorithm
acting on that object.

The present paper keeps the expansions.  What
Example~\ref{ex:p0-distributive-charge-transport} shows is that the expansions
which distributivity relates are distinguishable objects: they carry different
one-hole words, different intermediate values, and different charges, while
sharing one operator and one endpoint.  This is the sense in which AEG studies
process rather than result.  The value of an expression is one coordinate of
its expansions, not the object itself; the geometry developed in the following
sections is the geometry of the expansions.

\subsection{Point propagation and condition propagation}

The same one-hole word supports two elementary readings.  Given maps
$F_1,\ldots,F_n$ and a point $z_0$, the prefix operators
\[
  G_k=F_k\circ\cdots\circ F_1
\]
push the point forward:
\[
  z_k=G_k(z_0).
\]
This is the \emph{point propagation} used for paths.

Alternatively, given a family of output level sets $\{\Sigma_t\}$, the same
prefix operator pulls each level set back:
\[
  \Sigma_{k,t}=G_k^{-1}(\Sigma_t).
\]
This is \emph{condition propagation}.  It is contravariant: an outer condition
is transported toward the input by inverse image.  Sections~\ref{sec:p0-aes}
and~\ref{sec:p0-ripples} show that the first reading yields a path in a space of
arithmetic motions, while the second yields what we shall call ripple geometry.
